\SourceLeaf{322}%
\par\noindent\hspace*{\parindent} First, combining equations (1), (2) and (3) shows that \(qq+\alpha\alpha ll-uu-ss\) is a quantity of fourth order. Hence in place of (2) one can also write
\[
l=\frac{s\mathord{.}\sin U}{\alpha\mathord{.}\cos Q}(1-\tfrac1{24}qq+\tfrac1{24}uu)
\]
or, using (8), (9) and (10) to write
\[
s=\frac{r\cos\theta^2}{a\cos\varphi},\qquad \alpha\cos Q=\frac{\cos\theta\cos P}{\cos\varphi}
\]
one obtains
\[
l=\frac{r\cos\theta\sin U}{a\cos P}(1-\tfrac1{24}qq+\tfrac1{24}uu)
\]
Next, with equation (4) and a simple calculation, \(\frac{\cos\theta}{\cos P}=\frac{\BareRoot{(1-ee\sin P^2)}}{\cos P}\) expands to
\[
\frac{\cos\theta}{\cos P}=\frac{\BareRoot{(1-ee\sin B^2)}}{\cos B}(1+\frac{3ee\sin P^2}{8(1-ee\sin P^2)}\mathord{.}qq)
\]
which is accurate up to but excluding fourth order. Therefore, if in accordance with equation (6) \(T\) is also written in place of \(U\), we have
\[
l=\frac{r\BareRoot{(1-ee\sin B^2)}\mathord{.}\sin T}{a\cos B}(1-\frac{1-10ee\sin P^2}{24(1-ee\sin P^2)}\mathord{.}qq+\tfrac1{24}uu)
\]
\Needspace{14\baselineskip}
In the parenthesized part, one can now substitute \(\frac{\cos\varphi}{\BareRoot{(1-ee\sin P^2)}}\mathord{.}b\) for \(q\), then \(t\) for \(u\) and finally \(B\) for \(P\). As equations (5), (7) and (4) readily show, all these substitutions preserve the degree of accuracy. If we also write, for brevity,
\[
\BareRoot{(1-ee\sin B^2)}=k
\]
then we obtain (I)
\[
l=\frac{kr\sin T}{a\cos B}(1-\frac{(1-10ee\sin B^2)\cos\varphi^2}{24k^4}\mathord{.}bb+\tfrac1{24}tt)
\]

\Article{24}
Similarly, equation (1) becomes
\[
q=s\cos U(1-\tfrac1{12}qq+\tfrac1{12}ss+\tfrac18uu)
\]
and hence equation (5) becomes
